The principal affine $\mathfrak{sl}_2$ constraint admits a closed Virasoro field whose embedding preserves every collision, including the vacuum term.
The ghost relation determines the direction of this map.

\section{Currents, ghosts, and the differential}

All vector spaces are over $\mathbb C$.
Fix $k\in\mathbb C$ and put $t=k+2$.
Let $V^k(\mathfrak{sl}_2)$ be the universal affine vacuum vertex algebra with even generators $e,h,f$ and invariant form
\[
 (e,f)=1,\qquad (h,h)=2.
\]
The nonzero defining $\lambda$-brackets are
\begin{equation}\label{eq:affine-brackets}
 [h_\lambda e]=2e,\qquad [h_\lambda f]=-2f,
 \qquad [e_\lambda f]=h+k\lambda\mathbf1,
 \qquad [h_\lambda h]=2k\lambda\mathbf1.
\end{equation}
The other brackets follow by skew symmetry.
The affine algebra is universal, including at special levels.

Let $F_{bc}$ be the charged free-fermion vertex superalgebra with odd generators $b,c$ and
\[
 [b_\lambda c]=[c_\lambda b]=\mathbf1,
 \qquad [b_\lambda b]=[c_\lambda c]=0.
\]
Give $b,c$ cohomological degrees $-1,1$ and the affine generators degree zero.
Set $C_k=V^k(\mathfrak{sl}_2)\otimes F_{bc}$, with tensor-product vacuum $\mathbf1$ and translation $\partial$.
For every state $a$, use
\[
 Y(a,z)=\sum_{n\in\mathbb Z}a_{(n)}z^{-n-1},\qquad
 [a_\lambda d]=\sum_{n\geq0}\frac{\lambda^n}{n!}a_{(n)}d,
 \qquad :ad:=a_{(-1)}d.
\]
Thus $\{b_{(m)},c_{(n)}\}=\delta_{m+n,-1}$, and all five generating fields have $a_{(n)}\mathbf1=0$ for $n\geq0$.
The carrier consists of finite linear combinations of PBW monomials of negative modes.
It has no completion in ghost number, mode number, or tensor length.

Define the odd state and its residue by
\begin{equation}\label{eq:ds-charge}
 q=:(e-\mathbf1)c:,
 \qquad D=q_{(0)}=\Res_{z=0}Y(q,z)\,dz.
\end{equation}
The residue of $z^{-1}dz$ is $1$.
The defining operator products give
\begin{equation}\label{eq:ds-differential}
 Db=e-\mathbf1,\qquad Dh=-2:ec:,
 \qquad Df=:hc:+k\partial c,
 \qquad De=Dc=0.
\end{equation}
There are no singular contractions in $Y(q,z)Y(q,w)$.
Hence $[q_\lambda q]=0$, and the zero-mode commutator formula gives $D^2=0$.
The same formula gives, for every integer $n$ and homogeneous $a,d$,
\begin{equation}\label{eq:ds-derivation}
 D(a_{(n)}d)=(Da)_{(n)}d+(-1)^{|a|}a_{(n)}Dd,
 \qquad D\partial=\partial D.
\end{equation}
This is a differential vertex superalgebra at every level, including $k=-2$.
The noncritical hypothesis enters the conformal field below.

\section{The conformal field}

Assume $t\neq0$ and define
\begin{align}
 S&=\frac{\frac12:hh:+:ef:+:fe:}{2t},
 &G&=:(\partial b)c:,
 &T&=S+\frac12\partial h+G.\label{eq:ds-stress}
\end{align}
The conformal modes are $L_n=T_{(n+1)}$.

\begin{proposition}\label{prop:closed-stress}
The state $T$ is closed and is a conformal vector on $C_k$.
Its central charge and generating-field weights are
\begin{equation}\label{eq:ds-central}
 C(k)=\frac{3k}{k+2}-6k-2
      =13-6(t+t^{-1}),
 \qquad (\Delta_e,\Delta_h,\Delta_f,\Delta_b,\Delta_c)=(0,1,2,0,1).
\end{equation}
The field $q$ is primary of weight one.
In particular, $D$ commutes with every $L_n$.
\end{proposition}

\begin{proof}
The noncommutative Wick identity is
\[
 [a_\lambda:uv:]=:[a_\lambda u]v:
 +(-1)^{p(a)p(u)}:u[a_\lambda v]:
 +\int_0^\lambda [[a_\lambda u]_\mu v]\,d\mu,
\]
where $p$ denotes parity.
Applied to~\eqref{eq:affine-brackets}, it gives
\[
 [S_\lambda x]=(\partial+\lambda)x\quad(x=e,h,f),
 \qquad [S_\lambda S]=(\partial+2\lambda)S+\frac{k}{4t}\lambda^3\mathbf1.
\]
These are the Sugawara identities with the displayed invariant form.
In the ghost factor, the same identity gives
\[
 [G_\lambda b]=\partial b,
 \qquad [G_\lambda c]=(\partial+\lambda)c,
 \qquad [G_\lambda G]=(\partial+2\lambda)G-\frac16\lambda^3\mathbf1.
\]
Thus $G$ has central charge $-2$ and the stated weights.
The affine and ghost factors commute.
The mixed terms between $S$ and $\frac12\partial h$ supply
$(\partial+2\lambda)\frac12\partial h$, and
\[
 \left[\frac12\partial h{}_\lambda\frac12\partial h\right]
 =-\frac{k}{2}\lambda^3\mathbf1.
\]
Consequently
\begin{equation}\label{eq:TT}
 [T_\lambda T]=(\partial+2\lambda)T+\frac{C(k)}{12}\lambda^3\mathbf1.
\end{equation}
The zero mode of $\partial h$ vanishes, so $T_{(0)}=\partial$.
The weight statements follow from
\begin{align*}
 [T_\lambda e]&=\partial e,&
 [T_\lambda h]&=(\partial+\lambda)h-k\lambda^2\mathbf1,&
 [T_\lambda f]&=(\partial+2\lambda)f,\\
 [T_\lambda b]&=\partial b,&
 [T_\lambda c]&=(\partial+\lambda)c.
\end{align*}
The extra term in the $h$ bracket is a third-order pole; weight one does not make $h$ primary.

Directly from~\eqref{eq:ds-charge},
\[
 DS=:e\partial c:,
 \qquad D\left(\frac12\partial h\right)=-\partial:ec:,
 \qquad DG=:(\partial e)c:.
\]
Their sum is zero.
Since $e-\mathbf1$ is primary of weight zero, $c$ is primary of weight one, and their mutual bracket vanishes, the Wick identity gives
$[T_\lambda q]=(\partial+\lambda)q$.
Equivalently, the conformal modes commute with its residue $D$.
\end{proof}

The affine ghosts in~\eqref{eq:ds-charge} have weights $(0,1)$.
The reparametrization ghosts of weights $(2,-1)$ define a different complex.
In particular, the identity $D^2=0$ above imposes no equation $C(k)=26$.

\section{A reduced complex and its cohomology}

Put
\[
 j=:bc:,\qquad H=h+2j,\qquad F=f,\qquad E=e-\mathbf1,
 \qquad a=\frac{t-1}{2}.
\]
The elementary ghost identities needed below are
\begin{equation}\label{eq:ghost-identities}
 [j_\lambda j]=\lambda\mathbf1,
 \qquad [j_\lambda c]=-c,
 \qquad :jc:=-\partial c,
 \qquad :jj:=:(\partial b)c:-:b\partial c:.
\end{equation}
They follow by one or two contractions in the defining Clifford algebra.
Let $C_k^-$ be the vertex subalgebra generated by $H,F,c$, and let $C_k^+$ be the vertex subalgebra generated by $E,b$.
Their relevant brackets and differentials are
\begin{align}
 [H_\lambda H]&=2t\lambda\mathbf1,&
 [H_\lambda F]&=-2F,&
 [H_\lambda c]&=-2c,\label{eq:reduced-brackets}\\
 DH&=-2c,&
 DF&=:Hc:+t\partial c,&
 Dc&=0,\label{eq:reduced-differential}\\
 DE&=0,&Db&=E.\notag
\end{align}
All other brackets among $F,c$ vanish; all brackets in $C_k^+$ vanish.
Both subalgebras are subcomplexes.
Define
\begin{equation}\label{eq:W-and-L}
 W=F+\frac14:HH:+a\partial H,
 \qquad L=t^{-1}W.
\end{equation}

\begin{lemma}\label{lem:reduced-virasoro}
The state $W$ is closed.
The state $L$ satisfies the Virasoro bracket of central charge $C(k)$, and
\begin{equation}\label{eq:stress-homology}
 T-L=t^{-1}D(:bf:).
\end{equation}
\end{lemma}

\begin{proof}
Quasi-commutativity gives $:cH:=:Hc:+2\partial c$.
Therefore
\[
 D\left(\frac14:HH:\right)=-:Hc:-\partial c,
 \qquad D(a\partial H)=-2a\partial c.
\]
Equation~\eqref{eq:reduced-differential} and $2a=t-1$ prove $DW=0$.

For the Virasoro bracket, set $P=\frac14:HH:+a\partial H$.
The Heisenberg calculation using $[H_\lambda H]=2t\lambda$ gives
\[
 [P_\lambda P]=t(\partial+2\lambda)P
 +\left(\frac{t^2}{12}-2ta^2\right)\lambda^3\mathbf1.
\]
The Wick identity and $[F_\lambda H]=2F$ give
\[
 [F_\lambda P]=:HF:+t(\partial+\lambda)F,
 \qquad [P_\lambda F]= -:HF:+t\lambda F.
\]
Since $[F_\lambda F]=0$, division by $t^2$ yields
\[
 [L_\lambda L]=(\partial+2\lambda)L
 +\frac{1-24a^2/t}{12}\lambda^3\mathbf1.
\]
Here $1-24a^2/t=13-6(t+t^{-1})$.

Finally, $:fe:=:ef:-\partial h$ and~\eqref{eq:ghost-identities} give
\begin{align*}
 t(T-L)
 &=:(e-\mathbf1)f:-:hbc:-k:b\partial c:\\
 &=D(:bf:).
\end{align*}
All products in the middle expression use the affine--ghost tensor factors, so moving $h$ or $f$ through a ghost introduces no contraction.
\end{proof}

\begin{theorem}\label{thm:affine-comparison}
For every $k\neq-2$, there is an injective quasi-isomorphism of unital differential vertex superalgebras
\begin{equation}\label{eq:strict-map}
 \iota_T:(\Vir^{C(k)},0)\longrightarrow(C_k,D),
 \qquad \iota_T(\omega)=T,
\end{equation}
where $\Vir^{C(k)}$ is the universal Virasoro vacuum vertex algebra and $\omega$ its conformal generator.
The map is conformal for~\eqref{eq:ds-stress}.
In particular,
\[
 H^r(C_k,D)=0\quad(r\neq0),
 \qquad H^0(C_k,D)\simeq\Vir^{C(k)}.
\]
\end{theorem}

\begin{proof}
The bracket~\eqref{eq:TT} and $DT=0$ construct the unital map~\eqref{eq:strict-map} by the universal Virasoro presentation.

Normally ordered multiplication defines a linear map of complexes
\begin{equation}\label{eq:pbw-factorization}
 C_k^-\otimes C_k^+\longrightarrow C_k,
 \qquad u\otimes v\longmapsto u_{(-1)}v.
\end{equation}
This map is a linear isomorphism by the following PBW argument.
Give derivatives of $e,b,h,c,f$ filtration degrees $0,0,1,1,2$, respectively.
The affine and ghost brackets lower the sum of these degrees by at least one.
The associated graded algebra is therefore the polynomial algebra in the even jets and the exterior algebra in the odd jets.
The change of generators
\[
 (e,h,f,b,c)\longmapsto(E,H,F,b,c)
       =(e-1,h+2bc,f,b,c)
\]
and all its derivatives is an invertible change in this commutative algebra.
Thus the associated graded of~\eqref{eq:pbw-factorization} is an isomorphism.
Induction on the nonnegative filtration degree proves the assertion on the original vector spaces.
Each state has finite filtration degree.
Equation~\eqref{eq:ds-derivation} proves that the map is a map of complexes, using the tensor sign $(-1)^{|u|}$ on the second differential.
It is not asserted to be a tensor-product decomposition of vertex algebras.

Write $E_r=\partial^rE$ and $b_r=\partial^rb$ for $r\geq0$ in the commutative vertex algebra $C_k^+$.
Then
\[
 C_k^+=\mathbb C[E_r:r\geq0]\otimes\Lambda(b_r:r\geq0),
 \qquad D=\sum_{r\geq0}E_r\frac{\partial}{\partial b_r}.
\]
The odd operator $s=\sum_{r\geq0}b_r\partial/\partial E_r$ satisfies
$Ds+sD=N$, where $N$ counts the total number of $E_r,b_r$ factors.
Every polynomial contains finitely many factors and variables.
On its positive-$N$ subspace, $s/N$ is a contraction.
Thus $\mathbb C\mathbf1\to C_k^+$ is a quasi-isomorphism and the inclusion $C_k^-\to C_k$ is a quasi-isomorphism by~\eqref{eq:pbw-factorization}.

For $C_k^-$, use the nonnegative filtration in which each derivative of $H,c$ has degree one and each derivative of $F$ degree two.
Its associated graded differential is
\[
 dH_r=-2c_r,\qquad dF_r=\partial^r(H_0c_0),\qquad dc_r=0.
\]
Introduce even variables $w_r=\partial^r(F_0+H_0^2/4)$.
The map $(F_r,H_r,c_r)\mapsto(w_r,H_r,c_r)$ is polynomial and invertible, and $dw_r=0$.
Consequently the associated graded complex is
\begin{equation}\label{eq:graded-reduction}
 \mathbb C[w_r:r\geq0]\otimes
 \left(\mathbb C[H_r:r\geq0]\otimes\Lambda(c_r:r\geq0),
       -2\sum_{r\geq0}c_r\frac{\partial}{\partial H_r}\right).
\end{equation}
On the second factor, $s_-=-\frac12\sum_rH_r\partial/\partial c_r$ gives
$ds_-+s_-d=N_-$, the number of $H_r,c_r$ factors.
This proves that its cohomology consists of constants.

The Virasoro relation for $L$ constructs a map $\iota_L:\Vir^{C(k)}\to C_k^-$.
Filter its source by assigning degree two to all derivatives of $\omega$.
The PBW basis of the universal Virasoro vacuum module consists of the ordered products of $L_{-n}$, $n\geq2$, acting on the vacuum.
Its associated graded algebra is $\mathbb C[\partial^r\omega:r\geq0]$.
The associated graded of $\iota_L$ sends $\partial^r\omega$ to $w_r/t$ and is a quasi-isomorphism by~\eqref{eq:graded-reduction}.
The associated graded of its mapping cone is acyclic.
For a cone cycle, solve its boundary equation in its highest filtration quotient and subtract that boundary.
Induction terminates because the filtration is bounded below.
Hence $\iota_L$ is a quasi-isomorphism.

Equation~\eqref{eq:stress-homology} identifies $[T]=[L]$ in $H^0(C_k)$.
The two induced vertex homomorphisms agree on the universal generator, hence agree everywhere.
Thus $\iota_T$ is a quasi-isomorphism as well.
Since its source has zero differential, an element in its kernel has zero cohomology class in the source and must be zero.
This proves injectivity.
\end{proof}

Theorem~\ref{thm:affine-comparison} concerns the universal vacuum carrier.
It agrees with the PBW and reduction statements in \cite{AffineDSArakawa026}*{Theorem~5.8, p.~32; \S5.3, p.~33; Example~5.22, p.~39}.
No vanishing statement on an unrestricted category of affine modules is used in the proof.
The sign of the character in~\eqref{eq:ds-charge} differs from the $e+1$ convention of \cite{AffineDSArakawa026}*{\S5.1}; the substitution $(e,f,b,c)\mapsto(-e,-f,-b,-c)$ identifies the two conventions.

\begin{proposition}[Direction of a strict comparison]\label{prop:no-forward-map}
There is no unital vertex superalgebra homomorphism from $C_k$ to a purely even vertex algebra.
In particular, $\iota_T$ has no strict unital vertex retraction $C_k\to\Vir^{C(k)}$.
\end{proposition}
\begin{proof}
Parity preservation sends $b,c$ to zero.
The identity $b_{(0)}c=\mathbf1$ would then send $\mathbf1$ to zero, contradicting the unit condition.
\end{proof}
The inverse to~\eqref{eq:strict-map} exists after inverting quasi-isomorphisms.
The ordinary projection from closed states to $H^0(C_k)$ is defined on $\ker D$, not on every state of $C_k$ as a vertex homomorphism.

\section{The unital collision map on an affine coordinate chart}

Let $X=\mathbb A^1_{\mathbb C}$ with coordinate $z$.
For a differential vertex superalgebra $(V,d)$, define a left $D_X$-module
\[
 \mathcal M_V=\mathcal O_X\otimes_{\mathbb C}V,
 \qquad \nabla_{\partial_z}(f\otimes v)=f'\otimes v+f\otimes\partial v,
\]
with internal differential $1\otimes d$.
Its corresponding right module is $\mathcal A_V=\omega_X\otimes\mathcal M_V$; explicitly,
\begin{equation}\label{eq:right-connection}
 (f\,dz\otimes v)\cdot\partial_z
 =-f'\,dz\otimes v-f\,dz\otimes\partial v.
\end{equation}
These modules use algebraic tensor products.

On $X^2$, use coordinates $z,w$ and put $x=z-w$.
Let $j_2:X^2\setminus\Delta\hookrightarrow X^2$ and $\Delta:X\hookrightarrow X^2$ be the open and diagonal embeddings.
In left-module notation, set
\[
 \mathcal P_V=\mathcal O_{X^2}(*\Delta)\otimes V\otimes V,
 \qquad
 \mathcal R_V=V\otimes\mathbb C[w]\otimes x^{-1}\mathbb C[x^{-1}].
\]
The second expression is the coordinate presentation of the module supported on the diagonal, with actions
\begin{equation}\label{eq:diagonal-connection}
 \partial_z=\partial_x,
 \qquad \partial_w=\partial_w^{\mathrm{base}}-\partial_x+\partial_V.
\end{equation}
Multiplication by $x$ is interpreted modulo regular terms.
Tensoring these left modules with $dz\wedge dw$ identifies their right-module versions with
$j_{2*}j_2^*(\mathcal A_V\boxtimes\mathcal A_V)$ and $\Delta_!\mathcal A_V$.
The residue orientation is $dz\wedge dw=dx\wedge dw$ and $\Res x^{-1}dx=1$.

Define the full binary collision morphism by
\begin{equation}\label{eq:full-collision}
 \mu_V\bigl(f(z,w)a\boxtimes d\bigr)
 =\operatorname{PP}_x\bigl(f(w+x,w)Y(a,x)d\bigr).
\end{equation}
Here $\operatorname{PP}_x$ discards nonnegative powers of $x$.
For fixed $a,d$ and a meromorphic coefficient $f$ of finite pole order, only finitely many modes contribute to the principal part.
In particular, this formula also uses negative vertex products when $f$ has a pole.
It keeps the vacuum terms of all nonnegative products.

\begin{lemma}\label{lem:collision-chain}
Formula~\eqref{eq:full-collision} is a morphism of differential $D_{X^2}$-modules.
Its right-module form is the unital chiral product on $\mathcal A_V$.
\end{lemma}
\begin{proof}
For $\partial_z$, the chain rule and $Y(\partial a,x)=\partial_xY(a,x)$ identify the source connection with the first action in~\eqref{eq:diagonal-connection}.
For $\partial_w$, use
\[
 \partial_VY(a,x)d-\partial_xY(a,x)d=Y(a,x)\partial_Vd.
\]
This gives the second action in~\eqref{eq:diagonal-connection}.
Taking principal parts preserves both equations.
The differential identity follows from~\eqref{eq:ds-derivation}, with $D$ replaced by $d$.
Skew symmetry and the three-input Jacobi identity are the corresponding vertex identities after passage to principal parts; right-module conversion supplies the ordinary antisymmetry sign, together with the parity sign.
The vacuum embedding is
$\omega_X\to\mathcal A_V$, $f\,dz\mapsto f\,dz\otimes\mathbf1$.
For example,
\[
 \operatorname{PP}_x\bigl(x^{-1}Y(\mathbf1,x)a\bigr)=x^{-1}a,
\]
which is the unit collision under the residue identification.
This construction is the coordinate expression of \cite{AffineDSFrenkel026}*{\S7, Theorem~7.1 and Corollary~7.2, pp.~875--34--35}.
\end{proof}

\begin{theorem}[Chiral comparison on an affine chart]\label{thm:chiral-chart}
For $k\neq-2$, formula
\begin{equation}\label{eq:sheaf-map}
 \mathcal I_X:\mathcal A_{\Vir^{C(k)}}\longrightarrow\mathcal A_{C_k},
 \qquad f\,dz\otimes v\longmapsto f\,dz\otimes\iota_T(v)
\end{equation}
defines a unital quasi-isomorphism of differential chiral algebras on $\mathbb A^1$.
It preserves the full binary collision morphism:
\begin{equation}\label{eq:collision-square}
 (\Delta_!\mathcal I_X)\mu_{\Vir}
 =\mu_{C_k}\,j_{2*}j_2^*(\mathcal I_X\boxtimes\mathcal I_X).
\end{equation}
Every finite iterated collision, with its meromorphic coefficient and its vacuum outputs, satisfies the corresponding intertwining identity.
\end{theorem}
\begin{proof}
The state map $\iota_T$ commutes with translation, the vacuum, and the internal differentials.
It therefore respects~\eqref{eq:right-connection} and the unit embedding.
For every $n\in\mathbb Z$,
\[
 \iota_T(a_{(n)}d)=\iota_T(a)_{(n)}\iota_T(d).
\]
Substitution into~\eqref{eq:full-collision} proves~\eqref{eq:collision-square} for every coefficient $f$, including all its regular and polar parts.
Induction on the number of binary operations proves the iterated identities.
The underlying sheaf map is a quasi-isomorphism because Theorem~\ref{thm:affine-comparison} is a quasi-isomorphism of complexes of vector spaces and tensoring with each $\mathcal O_X$-stalk is exact.
\end{proof}

The lowest Virasoro collision reads
\begin{equation}\label{eq:TT-ope}
 Y(T,x)T\sim\frac{C(k)}{2x^4}\mathbf1
       +\frac{2T}{x^2}+\frac{\partial T}{x}.
\end{equation}
Thus $T_{(3)}T=C(k)\mathbf1/2$, $T_{(2)}T=0$, $T_{(1)}T=2T$, and $T_{(0)}T=\partial T$.
For a regular coefficient $g(x)$, the residue is
\[
 \Res_{x=0}g(x)Y(T,x)T\,dx
 =\frac{C(k)}{12}g'''(0)\mathbf1+2g'(0)T+g(0)\partial T.
\]
At $k=-8/3$ or $-7/2$, taking $g=x^3$ gives $13\mathbf1$.
The independent ghost collision is $\Res Y(b,x)c\,dx=\mathbf1$.
Neither vacuum output can be removed by a linear projection while preserving the original product.

\section{The binary collision complex}

For any $V$ above, put
\[
 P_V=j_{2*}j_2^*(\mathcal A_V\boxtimes\mathcal A_V),
 \qquad R_V=\Delta_!\mathcal A_V.
\]
Their internal differentials are $d_P(a\boxtimes b)=da\boxtimes b+(-1)^{|a|}a\boxtimes db$ and $d_R$ induced by $d$.
Define the unital binary collision complex
\begin{equation}\label{eq:binary-cone}
 K_2(V)=P_V[1]\oplus R_V,
 \qquad d_{K_2}(sp,r)=(-s d_Pp,\,\mu_V(p)+d_Rr).
\end{equation}
The suspension convention is $|sp|=|p|-1$.
The equality $d_R\mu_V=\mu_Vd_P$ proves $d_{K_2}^2=0$.

\begin{corollary}\label{cor:binary-quasi}
The componentwise map
\[
 K_2(\Vir^{C(k)})\longrightarrow K_2(C_k)
\]
defined by $\mathcal I_X\boxtimes\mathcal I_X$ on $P[1]$ and $\Delta_!\mathcal I_X$ on $R$ is a quasi-isomorphism.
\end{corollary}
\begin{proof}
Equation~\eqref{eq:collision-square} proves the total chain identity, including the off-diagonal component $\mu$.
Tensor products of complexes over $\mathbb C$ preserve quasi-isomorphisms.
The localization $\mathcal O_{X^2}\to\mathcal O_{X^2}(*\Delta)$ is flat, so the map on $P$ is a quasi-isomorphism.
The coordinate presentation~\eqref{eq:diagonal-connection} identifies the map on $R$ with a direct sum of copies of the state map tensored with $\mathcal O_X$, so it is a quasi-isomorphism as well.
The cone of the total map has a two-step filtration with these two acyclic cones as quotients.
It is acyclic by the long exact sequence of cohomology.
\end{proof}

The complex~\eqref{eq:binary-cone} retains the unit and is explicitly defined on $X^2$.
Its comparison does not require an augmentation of $C_k$.
An identification with a reduced chiral bar defined using an augmentation kernel would require such an augmentation, which is excluded by $b_{(0)}c=\mathbf1$.
An infinite totalization over configuration spaces requires its own definition and convergence theorem; the two-step argument above contains no such limit.

\section{Descent for the Drinfeld--Sokolov conformal structure}

The conformal structure~\eqref{eq:ds-stress} also specifies coordinate changes.
Let $X$ now be a smooth complex algebraic curve, and let $\widehat X\to X$ be its torsor of formal coordinates.
The Lie algebra of $\Aut\mathbb C[[z]]$ acts on $C_k$ by
\begin{equation}\label{eq:aut-action}
 -z^{n+1}\partial_z\longmapsto L_n\qquad(n\geq0).
\end{equation}
The Virasoro central term vanishes on this subalgebra.
All PBW states have nonnegative integral $L_0$-weights, and $L_n$ lowers weight by $n>0$.
The orbit of a weight-$N$ state lies in the span of the finitely many words in $L_1,\ldots,L_N$ with total weight loss at most $N$.
The positive part therefore integrates to formal coordinate changes, and integral $L_0$ integrates to $\mathbb G_m$.
The action is locally finite despite the infinite-dimensional weight-zero subspace.

Form the associated sheaf
\[
 \mathcal M^{\mathrm{DS}}_{C_k,X}
     =\widehat X\mathbin{\times}_{\Aut\mathbb C[[z]]}C_k,
 \qquad
 \mathcal A^{\mathrm{DS}}_{C_k,X}
     =\omega_X\otimes\mathcal M^{\mathrm{DS}}_{C_k,X}.
\]
Here $(u g,v)\sim(u,g v)$ for the right coordinate torsor and the left action~\eqref{eq:aut-action}.
The translation operator extends~\eqref{eq:aut-action} to regular formal vector fields and gives the connection $d+\partial\,dz$ in a coordinate trivialization.
This is the associated-bundle connection described in \cite{AffineDSFrenkel026}*{\S4.2, p.~875--18}.
The sheaf is the union of its finite-rank associated subbundles as an $\mathcal O_X$-module; the connection can raise their filtration.

\begin{proposition}\label{prop:global-descent}
For every smooth complex algebraic curve $X$ and $k\neq-2$, the map $\iota_T$ descends to a unital quasi-isomorphism of differential chiral algebras
\[
 \mathcal A_{\Vir^{C(k)},X}
       \longrightarrow\mathcal A^{\mathrm{DS}}_{C_k,X}.
\]
The local binary collision comparison of Corollary~\ref{cor:binary-quasi} descends to $X^2$.
\end{proposition}
\begin{proof}
The conformal map $\iota_T$ intertwines all $L_n$ and hence their integrated coordinate actions.
Proposition~\ref{prop:closed-stress} says that $D$ commutes with the same action.
Thus both the sheaf map and its differential are defined on the associated sheaves.
Translation compatibility gives horizontality.
The coordinate covariance of the vertex operation, equivalently the chiral construction in Lemma~\ref{lem:collision-chain}, glues the collision identity on overlaps.
Its unit is preserved before descent.
Locally in an \'etale coordinate, the underlying complex is $\mathcal O_X\otimes C_k$, so Theorem~\ref{thm:affine-comparison} proves the quasi-isomorphism.
The same local argument applies to the two terms of~\eqref{eq:binary-cone}; quasi-isomorphism is local.
\end{proof}

The superscript $\mathrm{DS}$ specifies the coordinate action, which uses $T$.
For example, $e$ and $b$ have weight zero, $f$ has weight two, and $h$ has third-order conformal pole $-2k\mathbf1$.
These are not the three weight-one primary currents of the Sugawara coordinate action.
The proposition therefore does not identify this associated sheaf with a separately prescribed affine-current bundle using other coordinate transitions.
Nor does it construct compatibility with clutching at nodal curves, a completed chiral bar, or a second semi-infinite differential.
Each such comparison has an additional carrier and an additional intertwining equation.
